%============================================================================%
% Philippines: rainfall & temperature mortality response functions
%
% Figures/tables are produced by adhoc/phl/01_phl_response_functions.R
% Compile from this directory:  latexmk -pdf phl_rainfall_mortality.tex
%============================================================================%
\documentclass[10pt, aspectratio=169]{beamer}

% dove is a beamer COLOR theme (greyscale, minimal), not a presentation theme
\usetheme{default}
\usecolortheme{dove}
\setbeamertemplate{navigation symbols}{}
\setbeamertemplate{caption}{\insertcaption}
\setbeamerfont{title}{series=\bfseries}
\setbeamerfont{frametitle}{series=\bfseries}
\setbeamercolor{frametitle}{fg=black}

\usepackage{booktabs}
\usepackage{graphicx}
\usepackage{amsmath}
\usepackage[T1]{fontenc}
\usepackage{lmodern}   % provides a bold sans face, so bold frame titles resolve

\graphicspath{{../figures/phl/}}

% panel-structure macros (\phlNobs, \phlNadmthree, ...) written by the R script
\input{../tables/phl/panel_facts.tex}

% section divider slides
\AtBeginSection[]{%
  \begin{frame}[plain]
    \vfill
    \centering
    {\usebeamerfont{title}\insertsectionhead}
    \vfill
  \end{frame}
}

\title{Rainfall, Temperature and Mortality in the Philippines}
\subtitle{Municipality-level response functions, \phlYearMin--\phlYearMax}
\date{\today}

\begin{document}

\begin{frame}[plain]
  \titlepage
\end{frame}

%============================================================================%
\section{Data}
%============================================================================%

\begin{frame}{The panel}
  \begin{itemize}
    \item New Philippines mortality panel, merged to gridded climate and
          GHSL settlement data.
    \item \textbf{\phlNadmthree} municipalities (adm3), nested in
          \textbf{\phlNadmtwo} provinces (adm2) and \textbf{\phlNadmone}
          regions (adm1).
    \item Annual, \phlYearMin--\phlYearMax, balanced. The pandemic years are
          retained and handled with a control rather than dropped (see the
          COVID section).
    \item Three age groups: 0--4, 5--64, 65+. The all-ages total
          (\texttt{agegroup} 0) is dropped.
    \item \textbf{\phlNobs} municipality $\times$ year $\times$ age
          observations in the estimation sample. BARMM is almost entirely
          absent (8 of 127 municipalities) --- see the maps.
  \end{itemize}

  \vspace{0.5em}
  \textbf{Climate.} Daily ERA5-Land precipitation and temperature, aggregated
  to annual sums of daily polynomial terms. CHIRPS precipitation is used as an
  alternative product.
\end{frame}

\begin{frame}{Two notes on the climate variables}
  \begin{itemize}
    \item The regressors are \emph{annual sums of daily powers},
          $\widetilde{P}^{(k)}_{it} = \sum_{d \in t} P_{id}^{k}$, not powers of
          the annual total. Response functions are therefore interpreted over
          \textbf{daily} values.
    \item Precipitation curves are evaluated over 0--300 mm/day and
          temperature over 15--35$^\circ$C, which spans the support observed in
          the Philippines.
  \end{itemize}
\end{frame}

\begin{frame}{The climate}
  \begin{center}
    \includegraphics[width=0.95\textwidth]{map_climate.png}
  \end{center}
  {\scriptsize Municipality means over 2006--2024, ERA5-Land. The eastern
  seaboard --- the typhoon-facing coast --- is both the wettest and, in
  eastern Mindanao, the most variable year to year. Grey: municipalities absent
  from the mortality panel --- almost all of BARMM (119 of its 127
  municipalities; only 8 appear in the panel), reflecting civil-registration
  coverage there rather than a merge failure.}
\end{frame}

\begin{frame}{The heterogeneity covariates}
  \begin{center}
    \includegraphics[width=0.95\textwidth]{map_covariates.png}
  \end{center}
  {\scriptsize Urban share is the GHSL population-based measure averaged over
  the sample; settlement class is each municipality's modal GHSL SMOD class;
  GDP per capita varies only at the region (adm1) level --- the visual
  reminder of why the income margin is coarse.}
\end{frame}

\begin{frame}{Mortality geography}
  \begin{center}
    \includegraphics[width=0.95\textwidth]{map_mortality.png}
  \end{center}
  {\scriptsize Mean winsorised death rates by municipality. Note the
  independent colour scales: 65+ rates are roughly twenty times the other
  groups. The sparse-looking north-Luzon interior for 0--4 reflects small
  cohorts and many zero-death municipality-years, the same cells the count
  Poisson and log-OLS estimators treat differently.}
\end{frame}

%============================================================================%
\section{Summary Statistics}
%============================================================================%

\begin{frame}{Summary statistics}
  \begin{center}
    \footnotesize
    \input{../tables/phl/summary_stats.tex}
  \end{center}
  \vspace{0.3em}
  {\scriptsize Estimation sample. Mortality rates are deaths per 100,000.
  Climate variables are annual totals except temperature, reported as the
  implied daily mean.}
\end{frame}

%============================================================================%
\section{Empirical Strategy}
%============================================================================%

\begin{frame}{Estimating equation}
  For municipality $i$, age group $a$, year $t$:
  \begin{align*}
    D_{ait} = \exp\Big(
      &\sum_{k=1}^{K_T} \beta^{T}_{ak}\, \widetilde{T}^{(k)}_{it}
      + \sum_{k=1}^{K_P} \beta^{P}_{ak}\, \widetilde{P}^{(k)}_{it}
      + \mu_{a i c(t)} + \nu_{a t r(i)} \Big) \cdot \varepsilon_{ait}
  \end{align*}
  where
  \begin{itemize}
    \item $D_{ait}$ is the death rate per 100,000, winsorised at the 99th
          percentile within age group;
    \item $\widetilde{X}^{(k)}_{it} = \sum_{d \in t} X_{id}^{k}$ is the annual
          sum of daily values raised to the $k$th power;
    \item $\mu_{a i c(t)}$ is a municipality $\times$ age $\times$
          pandemic-era effect, where $c(t)$ is the era containing year $t$
          (pre-2020, 2020--22, 2023--24);
    \item $\nu_{a t r(i)}$ is a province $\times$ year $\times$ age effect,
          where $r(i)$ is the province containing municipality $i$.
  \end{itemize}
  \vspace{0.3em}
  Default uses $K_P = K_T = 3$ and specification (4) below. Estimated by
  Poisson pseudo-maximum likelihood, weighted by
  population share, with standard errors clustered at the province (adm2)
  level.
\end{frame}

\begin{frame}{Response functions}
  The reported curve is the change in the death rate from a day at $x$ relative
  to a reference day $x_{\text{ref}}$:
  \[
    f^{P}_{a}(p) \;=\; \sum_{k=1}^{K_P} \hat\beta^{P}_{ak}
      \left( p^{k} - p_{\text{ref}}^{k} \right),
    \qquad p_{\text{ref}} = 0 \text{ mm}
  \]
  \[
    f^{T}_{a}(\tau) \;=\; \sum_{k=1}^{K_T} \hat\beta^{T}_{ak}
      \left( \tau^{k} - \tau_{\text{ref}}^{k} \right),
    \qquad \tau_{\text{ref}} = 27^{\circ}\mathrm{C}
  \]
  Confidence intervals are delta-method, using the clustered variance matrix.

  \vspace{0.8em}
  \textbf{Heterogeneity.} Each margin replaces $\beta^{P}_{ak}$ with
  $\beta^{P}_{akg}$, where $g = g(i,t)$ indexes the heterogeneity cell, so the
  precipitation polynomial is fully interacted with age $\times$ group:
  \[
    \sum_{k=1}^{K_P} \sum_{g} \beta^{P}_{akg}\,
      \mathbf{1}\{g(i,t) = g\}\, \widetilde{P}^{(k)}_{it}
  \]
\end{frame}

\begin{frame}{Fixed-effects specifications}
  \begin{center}
    \small
    \begin{tabular}{ll}
      \toprule
      & Fixed effects \\
      \midrule
      (1) Year $\times$ age      & $\mu_{ai} + \nu_{at}$ \\
      (2) + province-age trends  & $\mu_{ai} + \nu_{at} + \delta_{a r(i)} \cdot t$ \\
      (3) Saturated              & $\mu_{ai} + \nu_{a t r(i)}$ \\
      \textbf{(4) + pandemic control} & $\mu_{a i c(t)} + \nu_{a t r(i)}$ \\
      \bottomrule
    \end{tabular}
  \end{center}
  \vspace{0.5em}
  $\mu_{ai}$ is a municipality $\times$ age effect; $r(i)$ denotes the
  \textbf{province} (adm2) containing municipality $i$ and $c(t)$ the pandemic
  era containing year $t$. \textbf{(4) is the default} used throughout. These are the
  single-country analogues of the specifications used in the global
  Carleton-style regressions elsewhere in the project: with one country,
  country $\times$ year $\times$ age collapses to year $\times$ age.
\end{frame}

%============================================================================%
\section{Baseline Response Functions}
%============================================================================%

\begin{frame}{Precipitation response, by age group}
  \begin{center}
    \includegraphics[width=0.78\textwidth]{baseline_precip.png}
  \end{center}
  {\scriptsize Default specification (4), full 2006--2024 panel with the
  pandemic control. Change in deaths per 100,000 relative to a day with no rainfall. Shaded bands are 95\% confidence intervals.}
\end{frame}

\begin{frame}{Fit diagnostics for the default specification}
  \begin{center}
    \includegraphics[width=0.9\textwidth]{fit_diagnostics.png}
  \end{center}
  {\scriptsize Non-lean refit of the default model; squared correlation
  between actual and fitted is $0.62$ / $0.82$ / $0.79$ for ages 0--4 / 5--64 /
  65+, almost all of it from the fixed effects rather than from climate. The
  right-hand panels use relative error $(y-\hat\mu)/\hat\mu$, which is
  scale-free and unaffected by the weighting; population bins are formed
  \emph{within} age group.
  \\[0.3em]
  \textbf{Top left}: mass on the 45$^\circ$ line; the band at zero for 0--4 is
  the zero-death cells. \textbf{Bottom left}: no funnel, so the Poisson
  variance assumption is not obviously violated; 65+ is most dispersed.
  \textbf{Top right}: the time drift is \emph{entirely} a small-municipality
  phenomenon --- large municipalities sit flat on zero throughout, while small
  ones move from about $-6\%$ in 2007 to $+10\%$ by 2019 before the
  pandemic-era effect resets them. Since the estimator is population-weighted,
  it fits the large units and lets the small ones drift; the direction is
  consistent with improving rural death registration over the period.
  \textbf{Bottom right}: median error is flat across the whole size range, and
  the IQR fans out sharply below $\sim$3,000 people --- ordinary small-count
  noise, not bias.}
\end{frame}

\begin{frame}{Residual geography}
  \begin{columns}[T]
    \begin{column}{0.5\textwidth}
      \begin{center}
        \includegraphics[width=0.9\textwidth]{map_residuals.png}
      \end{center}
    \end{column}
    \begin{column}{0.48\textwidth}
      \vspace{2em}
      Mean relative error $(y-\hat\mu)/\hat\mu$ by municipality, pooled over
      years and ages; the colour scale is clipped at the 99th percentile of
      $|$error$|$.
      \vspace{0.5em}

      The pattern is close to spatial noise --- no coherent typhoon-corridor
      or coastal signature, and no island-group block. That is reassuring for
      the fixed-effects design: what the model misses looks idiosyncratic
      rather than geographic.
      \vspace{0.5em}

      Grey municipalities are those absent from the panel, overwhelmingly
      BARMM.
    \end{column}
  \end{columns}
\end{frame}

\begin{frame}{Temperature response --- not identified in this panel}
  \begin{columns}[T]
    \begin{column}{0.52\textwidth}
      \begin{center}
        \includegraphics[width=\textwidth]{baseline_tavg.png}
      \end{center}
    \end{column}
    \begin{column}{0.46\textwidth}
      \vspace{1em}
      The fitted curves are wrong-signed somewhere at every age: hotter days
      imply \emph{lower} mortality for 5--64 (significantly) and 65+, while
      for 0--4 the heat slope is positive but \emph{cold} days imply
      significantly lower mortality than a 27$^\circ$C day.
      \vspace{0.4em}

      Not a solver problem (diagnostics converge in a handful of iterations;
      tighter tolerances move coefficients by $\sim 3\times10^{-12}$) --- an
      identification problem:
      \begin{itemize}
        \item the curves reshape substantially across fixed-effects sets
              (next slide);
        \item the polynomial terms are 0.99 collinear after the fixed effects.
      \end{itemize}
      Temperature is retained as a control, not interpreted.
    \end{column}
  \end{columns}
\end{frame}

\begin{frame}{Why: the fixed effects absorb the temperature signal}
  \begin{columns}[T]
    \begin{column}{0.46\textwidth}
      \vspace{1.2em}
      \begin{center}
        \scriptsize
        \input{../tables/phl/identification.tex}
      \end{center}
      \vspace{0.8em}
      {\scriptsize Computed under specification (1), whose year $\times$ age
      effects are the mildest time controls. Temperature in the Philippines
      is spatially smooth and dominated by national interannual (ENSO) swings,
      so year effects remove almost all of it: the identifying variation left in
      daily temperature is \textbf{0.127$^\circ$C}. A quadratic over 15--35$^\circ$C
      cannot be fitted off that. Rainfall retains roughly five times as much.}
    \end{column}
    \begin{column}{0.52\textwidth}
      \begin{center}
        \includegraphics[width=\textwidth]{temp_fe_instability.png}
      \end{center}
      {\scriptsize Temperature response by fixed-effects set ($K_T=3$). The
      heat slope flips sign across the sets --- at 33$^\circ$C the 65+ response
      is $-0.003$ under year $\times$ age effects but $+0.002$ with
      province-age trends --- with no stable shape. The precipitation response
      at 200 mm/day, by contrast, stays positive for every age group across all
      four specifications.}
    \end{column}
  \end{columns}
\end{frame}

%============================================================================%
\section{Robustness}
%============================================================================%

\begin{frame}{Precipitation polynomial order}
  \begin{center}
    \includegraphics[width=0.95\textwidth]{robust_poly_order.png}
  \end{center}
  {\scriptsize Temperature polynomial held at order 3 throughout. Order 1
  forces monotonicity in daily rainfall and is shown for reference; see the
  polynomial-selection results for formal comparison.}
\end{frame}

\begin{frame}{Which polynomial order? Formal selection}
  \begin{center}
    \includegraphics[width=0.92\textwidth]{poly_selection.png}
  \end{center}
  \vspace{-0.3em}
  {\scriptsize \textbf{The two criteria disagree, and the disagreement is the
  finding.} In-sample AIC/BIC fall monotonically in the order (best at 5) ---
  the usual in-sample reward for flexibility, which at $n=84{,}075$ the
  parameter penalty never offsets. But \textbf{out-of-sample deviance is
  flat}: the spread across orders 1--5 is $1.5\times10^{-6}$ against a
  fold-to-fold standard error of $2.9\times10^{-4}$, roughly \textbf{190 times
  larger}. None of the extra curvature generalises. Order 3 (the default) is
  defensible; so is order 1. Computed by
  \texttt{03\_phl\_poly\_selection.R}, which is kept out of the main pipeline
  because it refits the model 30 times.}
\end{frame}

\begin{frame}{Two caveats on the selection exercise}
  \begin{itemize}
    \item \textbf{The criteria shown are from an \emph{unweighted} refit.}
          Under the estimation default the likelihood is on the scale of
          population \emph{shares}, so the total quasi-log-likelihood is about
          $-751$ across 84k rows and AIC's $2k$ penalty swamps it: the gap
          between orders 1 and 2 is $5.85$ against a pure penalty of
          $2\times3 = 6$. Weighted AIC/BIC here are parameter-counting, not
          model comparison, and would trivially select order 1. Both variants
          are in \texttt{poly\_selection.csv}.
    \item \textbf{Folds are random at the observation level}, not blocked by
          municipality or year --- blocking would leave held-out rows with no
          identified fixed effect to predict from. So this measures functional
          form, not forecasting skill into new places or years. About 82\% of
          held-out rows are predictable; the rest fall in fixed-effect cells
          that vanish from the training fold.
  \end{itemize}
  \vspace{0.5em}
  {\scriptsize Taken with the wide confidence bands throughout, the reading is
  that this panel pins down the \emph{level} of the rainfall--mortality
  relationship far better than its \emph{shape}.}
\end{frame}

\begin{frame}{Fixed-effects specification}
  \begin{center}
    \includegraphics[width=\textwidth]{robust_fe.png}
  \end{center}
  {\scriptsize Specifications (1)--(4) as defined above.}
\end{frame}

\begin{frame}{Climate product: ERA5-Land vs CHIRPS}
  \begin{center}
    \includegraphics[width=0.85\textwidth]{robust_product.png}
  \end{center}
  {\scriptsize Both columns use ERA5-Land temperature controls; only the
  precipitation product differs.}
\end{frame}

\begin{frame}{Unit of analysis: municipality vs province}
  \begin{center}
    \includegraphics[width=0.85\textwidth]{robust_adm2.png}
  \end{center}
  {\scriptsize Left: the default municipality (adm3) panel. Right: collapsed to
  provinces (adm2), 6,270 observations, with fixed effects shifted up one level
  (province $\times$ age $\times$ era and region $\times$ year $\times$ age).
  Point estimates are close --- at 200 mm/day and 65+, $0.0039$ against
  $0.0026$ --- so the result is not an artefact of the spatial unit. Province
  aggregation actually \emph{retains} slightly more identifying rainfall
  variation after fixed effects (23.2\% vs 17.8\%), because the coarser
  fixed effects absorb less.}
\end{frame}

\begin{frame}{How the aggregation is done}
  Two choices matter when collapsing municipalities into provinces:
  \begin{itemize}
    \item \textbf{The outcome aggregates exactly.} A province death rate is
          total deaths over total population, which \emph{is} the
          population-weighted mean of its municipality rates --- no modelling
          choice. Winsorisation is re-applied at the province level, since the
          99th percentile of a different distribution is a different number.
    \item \textbf{The climate regressors are averaged as polynomial terms},
          not recomputed from averaged rainfall:
          \[
            \textstyle\sum_i w_i \big( \sum_d P_{id}^{k} \big)
            \quad\text{not}\quad
            \big( \sum_i w_i \sum_d P_{id} \big)^{k}
          \]
          Because the model is linear in these terms, the first makes the
          province response function exactly the population-weighted average
          of its municipalities' response functions. The second is a different
          object --- a function of average rainfall rather than the average of
          the function --- and would understate extremes by Jensen's
          inequality.
    \item Weights are \textbf{time-invariant} (period-mean municipal
          population), so that a mortality shock which moves population cannot
          also move measured climate exposure.
  \end{itemize}
\end{frame}

\begin{frame}{Estimator: Poisson (rate and count) vs OLS (levels and logs)}
  \begin{center}
    \includegraphics[width=\textwidth]{robust_estimator.png}
  \end{center}
  {\scriptsize Default specification (4). The default is Poisson on the
  winsorised rate, population-weighted; the count model is Poisson on raw
  deaths with a $\log$(population) offset, unweighted, and keeps zero-death
  cells. The first three panels are all in log points and directly comparable;
  OLS in levels (right) is in deaths per 100,000 --- note the free $y$ scales.
  The log-OLS sample drops zero-deathrate cells, mostly ages 0--4 ($\sim$5\% of
  that group).}
\end{frame}

\begin{frame}{Dynamics: lags, cumulative effect, and a lead placebo}
  \begin{center}
    \includegraphics[width=\textwidth]{dynamics.png}
  \end{center}
  {\scriptsize Default specification (4) with contemporaneous, one- and
  two-year-lagged precipitation polynomials entered jointly (first four
  panels; sample starts two years later to accommodate the lags).
  \textbf{No harvesting}: the lags enter weakly positive, not negative, and
  the cumulative response exceeds the contemporaneous one (65+ at 200 mm/day:
  $0.0059$ vs $0.0026$), though the cumulative CI spans zero. The rightmost
  panel is a falsification test from a separate regression; annual rainfall is
  serially uncorrelated within municipality (AR(1) $\approx -0.04$), so the
  lead cannot proxy current exposure. Its CI spans zero at every age, but the
  test is low-powered --- the placebo CI is wide enough to contain effects the
  size of the main estimate.}
\end{frame}

\begin{frame}{Leave-one-region-out}
  \begin{center}
    \includegraphics[width=0.9\textwidth]{loro.png}
  \end{center}
  {\scriptsize Response at 200 mm/day, re-estimated 17 times dropping one
  region (adm1) at a time. Dashed line marks the full-sample estimate. A
  result driven by a single region --- e.g. the typhoon corridor --- would
  show one point far from the dashed line.}
\end{frame}

\begin{frame}{Weights and winsorisation}
  \begin{columns}[T]
    \begin{column}{0.44\textwidth}
      \begin{center}
        \includegraphics[width=\textwidth]{robust_weights.png}
      \end{center}
    \end{column}
    \begin{column}{0.54\textwidth}
      \begin{center}
        \includegraphics[width=\textwidth]{robust_winsor.png}
      \end{center}
    \end{column}
  \end{columns}
  {\scriptsize Left: population-weighted (default) vs unweighted. The gap is
  real: at 200 mm/day the unweighted 65+ estimate roughly doubles ($0.0039 \to
  0.0072$) while 5--64 halves ($0.0031 \to 0.0017$), pointing to larger old-age
  effects in smaller municipalities --- consistent with the rural/urban
  heterogeneity section. Right: winsorisation at p99 (default, inherited from
  the Carleton pipeline), p95, and none --- estimates barely move.}
\end{frame}

\begin{frame}{Clustering}
  \begin{center}
    \small
    \input{../tables/phl/clustering.tex}
  \end{center}
  \vspace{0.4em}
  {\scriptsize Standard errors of the response at 200 mm/day under alternative
  clustering; point estimates identical by construction. Region-level
  clustering (17 clusters) barely moves the SEs but its asymptotics are poor.
  \textbf{Two-way province-and-year clustering raises the SEs by 40--100\%},
  and at 200 mm/day no age group then excludes zero --- serial or common
  shocks within years matter for inference here.}
\end{frame}

%============================================================================%
\section{COVID}
%============================================================================%

\begin{frame}{The pandemic sits inside the estimation window}
  \begin{center}
    \includegraphics[width=\textwidth]{mortality_trend.png}
  \end{center}
  {\scriptsize Population-weighted mean death rate by year. Mortality at 65+
  rises about 45\% in 2021 and does not return to its pre-pandemic level.
  Year $\times$ age effects absorb the national component of this, but not any
  covariance between local pandemic severity and local weather.}
\end{frame}

\begin{frame}{COVID sample sensitivity is modest under the saturated FEs}
  \begin{center}
    \includegraphics[width=\textwidth]{robust_covid.png}
  \end{center}
  {\scriptsize All four samples estimated under the saturated specification (3),
  so the comparison is not mediated by the pandemic control. The samples differ
  by well under one standard error --- at 200 mm/day and age 65+, $0.0059$,
  $0.0060$, $0.0054$ and $0.0069$ against a standard error of about $0.003$.
  The province $\times$ year $\times$ age effects already absorb most of the
  pandemic shock. Under the weaker year $\times$ age effects the same comparison
  showed the full sample roughly doubling the effect, so how much COVID matters
  is itself a function of the fixed effects.}
\end{frame}

\begin{frame}{The default: a pandemic control, not a truncated sample}
  \begin{center}
    \includegraphics[width=\textwidth]{pandemic_control.png}
  \end{center}
  \vspace{-0.3em}
  {\scriptsize The saturated specification already absorbs any pandemic shock
  common to a province $\times$ year $\times$ age cell, so the useful control is
  one level lower: a separate intercept for each \textbf{municipality $\times$
  age in each pandemic era} (pre-2020, 2020--22, 2023--24), which soaks up
  locality-specific COVID mortality while leaving within-era weather variation
  intact. At 200 mm/day for ages 5--64 it returns $0.0031$, against $0.0035$ on
  the pre-COVID sample and $0.0065$ with no control --- i.e. it recovers
  essentially the pre-COVID estimate while using all 19 years, with standard
  errors about 15\% tighter --- which is why this is the default here, rather
  than discarding five years of data.}
\end{frame}

%============================================================================%
\section{Heterogeneity}
%============================================================================%

\begin{frame}{Urban population share}
  \begin{center}
    \includegraphics[width=\textwidth]{het_urban.png}
  \end{center}
  {\scriptsize Precipitation polynomial interacted with age $\times$
  urban-share bin. Bins are \textbf{terciles of the urban-share
  distribution}, not equal-width: urban share is left-skewed (median 69\%), so
  equal-width bins would put 3.4 times as many municipality-years in the top
  bin as the bottom and estimate the rural response off far less data. Each
  bin here holds about 9,600 observations.}
\end{frame}

\begin{frame}{Urban population share --- five bins}
  \begin{center}
    \includegraphics[width=\textwidth]{het_urban5.png}
  \end{center}
  {\scriptsize Quintiles of the urban-share distribution, about 5,800
  municipality-years each. The finer grid buys little: no monotone gradient in
  urbanisation emerges, and only the most-urban bin for ages 0--4 excludes zero
  at 200 mm/day ($0.025$ $[0.007,\,0.042]$). Splitting further mostly costs
  precision.}
\end{frame}

\begin{frame}{Is the urban gradient real? Three tests}
  \begin{center}
    \includegraphics[width=0.82\textwidth]{urban_gradient.png}
  \end{center}
  \vspace{-0.4em}
  {\scriptsize Response at 200 mm/day by urban quintile. The eye is drawn to
  the top-right point; the tests below ask whether that is warranted.}
\end{frame}

\begin{frame}{Testing the threshold hypothesis directly}
  \begin{center}
    \includegraphics[width=0.8\textwidth]{urban_threshold.png}
  \end{center}
  \vspace{-0.3em}
  {\scriptsize Difference in the 200 mm/day response between municipalities
  above and below each urban-share cutoff. \textbf{No cutoff yields a
  significant difference at any age} --- the smallest $p$-value across all 18
  tests is $0.11$. Testing the top quintile against \emph{all} others pooled
  (the most powerful version of ``only the really urban places'') gives
  $+0.013$ $(p=0.13)$ for ages 0--4, $+0.0005$ $(p=0.88)$ for 5--64 and
  $+0.0002$ $(p=0.95)$ for 65+.}
\end{frame}

\begin{frame}{Splitting the top of the distribution finer}
  \begin{center}
    \includegraphics[width=\textwidth]{het_urban_topbins.png}
  \end{center}
  \vspace{-0.2em}
  {\scriptsize Bins chosen to be fine where the hypothesis lives:
  0--60, 60--80, 80--90, 90--95, 95--100\% urban. At 200 mm/day the
  significant cells are (80--90]\ and (90--95]\ for ages 0--4, (90--95]\ for
  65+ --- and \textbf{[0--60], the \emph{least} urban bin, for 65+}. The
  \textbf{most} urban bin (95--100\%) is not significant at any age.}
\end{frame}

\begin{frame}{What the urban result actually is}
  \begin{itemize}
    \item \textbf{The threshold hypothesis is not supported.} No cutoff from
          60\% to 95\% separates damaging from non-damaging municipalities,
          and the top quintile is indistinguishable from the rest.
    \item \textbf{Which bin is ``significant'' moves with the binning.} With
          quintiles it was (88.8--100]; with the finer grid it is (90--95],
          while (95--100] is not; and for 65+ the \emph{least} urban bin also
          clears. That instability is the signature of noise rather than a
          threshold.
    \item \textbf{But there is something --- and it is about age, not
          urbanisation.} For ages 0--4 the point estimate is positive in
          \emph{every} bin ($0.014$, $0.007$, $0.019$, $0.027$, $0.020$) and
          the pooled estimate sits just short of significance
          ($0.015$, s.e.\ $0.008$). The most defensible reading is a broad
          infant response to extreme rainfall across the whole urban
          distribution, which bin-by-bin testing makes \emph{look} like an
          urban threshold.
  \end{itemize}
  \vspace{0.2em}
  {\scriptsize Distinguishing these needs a mechanism, not more bins:
  cause-of-death for 0--4 (waterborne disease, drowning) is the test, and the
  data is already cleaned on the cluster.}
\end{frame}

\begin{frame}{Settlement type (GHSL degree of urbanisation)}
  \begin{center}
    \includegraphics[width=\textwidth]{het_smod.png}
  \end{center}
  {\scriptsize GHSL SMOD classes collapsed to rural (10--13), peri-urban
  (21--22) and urban centre (23, 30).}
\end{frame}

\begin{frame}{Income}
  \begin{center}
    \includegraphics[width=\textwidth]{het_income.png}
  \end{center}
  {\scriptsize Terciles of region-level average GDP per capita. Income varies
  only across the 17 regions, so this margin is coarse.}
\end{frame}

\begin{frame}{Long-run precipitation (adaptation margin)}
  \begin{center}
    \includegraphics[width=\textwidth]{het_lrprecip.png}
  \end{center}
  {\scriptsize Terciles of region-level long-run mean annual precipitation.}
\end{frame}

\begin{frame}{Coastal exposure}
  \begin{columns}[T]
    \begin{column}{0.30\textwidth}
      \begin{center}
        \includegraphics[width=\textwidth]{map_coastal.png}
      \end{center}
    \end{column}
    \begin{column}{0.68\textwidth}
      \begin{center}
        \includegraphics[width=\textwidth]{het_coastal.png}
      \end{center}
    \end{column}
  \end{columns}
  \vspace{-0.2em}
  {\scriptsize Exposure is the share of each municipality's perimeter facing
  the sea, computed from the adm3 shapefile by subtracting shared borders from
  the perimeter (\texttt{04\_phl\_coastal.R}). 41\% of municipalities are
  landlocked; the coastal ones are split at their median (37\%). Validation:
  the measure returns 0 for inland Ilocos Norte and 1 for Batanes and the
  Babuyan and Polillo islands.}
\end{frame}

\begin{frame}{Coastal exposure: the sign is the wrong way round}
  \begin{itemize}
    \item At 200 mm/day the \textbf{landlocked} bin carries the largest
          effects and is the \emph{only} one whose confidence intervals exclude
          zero --- ages 0--4: $0.025$ $[0.006,\,0.045]$; 65+: $0.0068$
          $[0.0005,\,0.0131]$. Highly-coastal municipalities show smaller,
          insignificant responses at every age.
    \item That is the opposite of a storm-surge story, and consistent with two
          non-exclusive readings: heavy rain drains quickly to the sea on the
          coast but \textbf{accumulates in inland river valleys and basins},
          and coastal municipalities facing chronic typhoon exposure are
          \textbf{better adapted} (warning systems, evacuation, building
          practice) than inland ones that see extreme rainfall rarely.
    \item \textbf{Do not over-read it.} Only 2 of 9 age $\times$ bin cells
          exclude zero, the intervals are wide, and the bins are cross-
          sectional so the contrast is between places rather than within them.
          A convincing version needs the daily extreme-rainfall measures and
          the cause-of-death split --- drowning and injury deaths should load
          on this margin if the flooding mechanism is right.
  \end{itemize}
\end{frame}

\begin{frame}{Marginal effects at 50, 100 and 200 mm/day}
  \begin{center}
    \includegraphics[width=0.9\textwidth]{marginal_effects.png}
  \end{center}
  {\scriptsize Point estimates with 95\% confidence intervals, read off the
  fitted response curves at the indicated daily rainfall levels.}
\end{frame}

%============================================================================%
\section{Alternative Functional Form}
%============================================================================%

\begin{frame}{Polynomial in annual, rather than daily, rainfall}
  \begin{columns}[T]
    \begin{column}{0.54\textwidth}
      \begin{center}
        \includegraphics[width=\textwidth]{annual_form.png}
      \end{center}
    \end{column}
    \begin{column}{0.44\textwidth}
      \vspace{1.5em}
      Everything else in this deck uses the Carleton form, where the order-$k$
      regressor is $\sum_d P_{id}^{k}$ and the curve is read over \emph{daily}
      rainfall. Here the regressor is $\big(\sum_d P_{id}\big)^{k}$ instead, so
      the curve is read over \emph{annual} rainfall.
      \vspace{0.6em}

      These answer different questions --- what one wet day does, versus what a
      wet year does --- and coincide only if rainfall were constant within the
      year.
      \vspace{0.6em}

      Under the default fixed effects the annual form is very imprecise: every
      confidence interval spans zero across 1,500--4,000 mm/year. Province
      $\times$ year effects leave little annual-total variation to work with.
    \end{column}
  \end{columns}
  \vspace{0.2em}
  {\scriptsize Dashed line marks the reference, the sample median annual
  rainfall. Default specification (4), full panel.}
\end{frame}

%============================================================================%
\section{Comparison with Global Estimates}
%============================================================================%

\begin{frame}{Philippines vs the project's global response}
  \begin{center}
    \includegraphics[width=0.95\textwidth]{global_overlay.png}
  \end{center}
  {\scriptsize The PHL response (default specification (4)) overlaid on the
  global precipitation response from the project's Carleton-sample regressions
  (ERA5-Land, main global specification: adm2 $\times$ age and country
  $\times$ year $\times$ age effects). Both are Poisson, so both curves are in
  log points. At 200 mm/day the PHL 5--64 response ($0.0031$) is roughly an
  eighth of the global estimate ($0.0246$); 65+ is comparable ($0.0039$ vs
  $0.0057$). Samples, admin units and fixed effects differ, so read this as
  context rather than a like-for-like test.}
\end{frame}

%============================================================================%
\section{Drought vs Extreme Rainfall (Prototype)}
%============================================================================%

\begin{frame}{Separating drought from extreme rainfall: design}
  Drought and flood mortality operate at different timescales, so they enter
  as \emph{separate exposures} in one regression:
  \[
    D_{ait} = \exp\Big( \underbrace{f_a(\text{daily polynomial})}_{\text{wet-day extremes}}
      + \underbrace{g_a(\text{annual anomaly bins})}_{\text{cumulative deficit/surplus}}
      + \text{temp} + \mu_{a i c(t)} + \nu_{a t r(i)} \Big)\cdot\varepsilon_{ait}
  \]
  \begin{itemize}
    \item Prototype drought index: within-municipality \textbf{z-score of
          annual rainfall} (a poor man's SPI --- monthly totals for a real SPI
          need the cluster climate build). Bins: severe ($z<-1.5$), moderate
          ($-1.5$ to $-0.5$), normal (reference), wet, very wet.
    \item \textbf{Identification checked first} (table below): the
          absorption is substantial --- the saturated effects remove
          three-quarters of the z-score's variation --- but far from fatal:
          the severe-drought \emph{dummy} keeps about half its variation,
          versus 9\% for temperature. PHL droughts have more subnational
          texture than the ENSO narrative suggests.
    \item Bin shares: severe drought 6.9\% of municipality-years, moderate
          25\%; symmetric on the wet side.
  \end{itemize}
  \begin{center}
    \scriptsize
    \input{../tables/phl/drought_absorption.tex}
  \end{center}
\end{frame}

\begin{frame}{Prototype results: annual-anomaly bins, by fixed effects}
  \begin{center}
    \includegraphics[width=\textwidth]{drought_prototype.png}
  \end{center}
  {\scriptsize Effect of a drought/wet year relative to a normal rainfall year
  ($|z| \le 0.5$), estimated alongside the daily precipitation polynomial and
  temperature controls. Under the weak year $\times$ age effects, drought
  years show \emph{higher} mortality (65+ moderate drought: $+5.5\%$), but
  this collapses --- and for 65+ turns significantly \emph{negative}
  ($-1.0\%$, se $0.4\%$) --- once province $\times$ year effects are added,
  the signature of regional confounding rather than a drought effect. The
  precise near-zero under the default is itself informative: in this wet,
  non-agrarian-margin country there is no detectable drought mortality. The
  multi-country panel, with drought-prone regions and staggered timing, is
  where this margin gets its real test.}
\end{frame}

%============================================================================%
\section{Caveats}
%============================================================================%

\begin{frame}{Caveats and next steps}
  \begin{itemize}
    \item \textbf{The pandemic control is doing real work only for 5--64.}
          It reproduces the pre-COVID estimate closely for ages 5--64, but for
          65+ it lands below both the uncontrolled and the pre-COVID estimate.
          Under the saturated fixed effects the COVID-sample sensitivity is in
          any case well within one standard error --- it was large only under
          the weaker year $\times$ age effects.
    \item \textbf{The temperature response is not identified.} Every age group
          is wrong-signed somewhere --- hotter days imply lower mortality for
          5--64 and 65+, colder days imply lower mortality for 0--4 --- and the
          shapes reshape substantially across fixed-effects sets and polynomial
          orders. After the fixed effects only 0.127$^\circ$C of daily
          temperature variation remains and the polynomial terms are 0.99
          collinear (VIF 55), so the curves are essentially extrapolation.
          Temperature is kept as a control and should not be interpreted.
    \item \textbf{Income variation is coarse.} \texttt{gdppc\_adm1\_avg} takes
          only seven distinct values across 17 regions, so the income margin is
          close to a regional fixed effect.
    \item \textbf{Clustering.} Standard errors are clustered at the province
          (adm2, \phlNadmtwo\ clusters); the default fixed effects are also at
          province $\times$ year $\times$ age, so the identifying variation is
          within-province. Two-way province-and-year clustering raises the SEs
          by 40--100\% and removes what significance there is at 200 mm/day
          --- the inference is fragile to common within-year shocks.
    \item \textbf{Coverage.} Roughly 3\% of municipality-years are missing
          ERA5-Land climate and drop out; CHIRPS has full coverage, which is
          worth keeping in mind when comparing the two products.
    \item \textbf{Interpretation.} Curves are conditional on the fixed effects
          and are not causal estimates of typhoon exposure; extreme-rainfall
          days are sparse, so the upper end of each curve (beyond roughly
          150 mm/day) is identified off very few observations and is
          correspondingly imprecise.
  \end{itemize}
\end{frame}

\end{document}
